% Part: computability
% Chapter: computability-theory
% Section: k-1

\documentclass[../../../include/open-logic-section]{subfiles}

\begin{document}

\olfileid{cmp}{thy}{k1}
\olsection{न्यूनीकरणक्षमतेचे एक उदाहरण}

\olref[ppr]{prop:reduce} चा एक उपयोग पाहू.

\begin{prop}
\ollabel{prop:k1}
पुढीलप्रमाणे घ्या:
\[
K_1 = \Setabs{e}{\cfind{e}(0) \fdefined}.
\]
मग $K_1$ संगणनक्षमपणे प्रगणनीय आहे, पण संगणनक्षम नाही.
\end{prop}

\begin{proof}
$K_1 = \Setabs{e}{\lexists[s][T(e,0,s)]}$ असल्याने,
\olref[eqc]{thm:exists-char} नुसार $K_1$ संगणनक्षमपणे
प्रगणनीय आहे.

$K_1$ संगणनक्षम नाही हे दाखवण्यासाठी $K_0$ चे
त्याच्याकडे न्यूनीकरण होते हे दाखवू.

\begin{explain}
हे थोडे गुंतागुंतीचे आहे, कारण $K_1$ वापरून आपण फक्त
विशिष्ट आदान $0$ पासून सुरू होणाऱ्या संगणनांबद्दल प्रश्न
विचारू शकतो. अशा प्रश्नांची उत्तरे देणारा हुशार मित्र
तुमच्याकडे आहे असे समजा (अशा मित्रांना ``ओरॅकल'' म्हणतात).
आता कोणीतरी तुम्हाला $\tuple{e,
  x}$ ही जोडी $K_0$ मध्ये आहे का, म्हणजे यंत्र $e$ हे
आदान $x$ वर थांबते का, असे विचारते. तुम्ही दुसरे यंत्र,
$e_x$, तयार करू शकता. ते \emph{कोणतेही} आदान मिळाले तरी
त्याकडे दुर्लक्ष करते आणि त्याऐवजी~$e$ हे यंत्र आदान
$x$ वर चालवते. मग यंत्र $e$ आदान $x$ वर थांबते का हा
प्रश्न, यंत्र $e_x$ आदान $0$ वर (किंवा इतर कोणत्याही
आदानावर) थांबते का या प्रश्नाशी समतुल्य आहे. आता नवीन
यंत्र $e_x$ आदान $0$ वर थांबते का हे मित्राला विचारा;
बदललेल्या प्रश्नाचे त्याचे उत्तर मूळ प्रश्नाचे उत्तरही देते.
यातून $K_0$ चे~$K_1$ कडे हवे ते न्यूनीकरण मिळते.
\end{explain}

सार्वत्रिक आंशिक संगणनक्षम फलन वापरून $f$ हे पुढीलप्रमाणे
व्याख्यायित त्रिस्थानी फलन असू द्या:
\[
f(x,y,z) \simeq \cfind{x}(y).
\]
$f$ आपल्या तिसऱ्या आदानाकडे पूर्णपणे दुर्लक्ष करते, हे
लक्षात घ्या. $f = \cfind{e}[3]$ होईल असा निर्देशांक $e$
निवडा; म्हणजे
\[
\cfind{e}[3](x,y,z) \simeq \cfind{x}(y).
\]
$s$-$m$-$n$ प्रमेयानुसार असे फलन $s(e,x,y)$ अस्तित्वात
आहे की प्रत्येक $z$ साठी
\begin{align*}
\cfind{s(e,x,y)}(z) & \simeq \cfind{e}[3](x,y,z) \\
& \simeq \cfind{x}(y).
\end{align*}

\begin{explain}
वरील अनौपचारिक युक्तिवादाच्या दृष्टीने $s(e,x,y)$ हा
अशा यंत्राचा निर्देशांक आहे की ते कोणतेही आदान $z$
मिळाले तरी त्याकडे दुर्लक्ष करते आणि $\cfind{x}(y)$
संगणित करते.
\end{explain}

विशेषतः, आपल्याला
\[
\cfind{s(e,x,y)}(0) \fdefined \quad \text{तेव्हाच} \quad
\cfind{x}(y) \fdefined.
\]
म्हणजे $\tuple{x, y} \in K_0$ असेल तेव्हाच $s(e,x,y) \in
K_1$. म्हणून पुढीलप्रमाणे व्याख्यायित फलन $g$:
\[
g(w) = s(e,(w)_0,(w)_1)
\]
हे $K_0$ चे~$K_1$ कडे न्यूनीकरण आहे.
\end{proof}

\end{document}
